\chapter{The inductive condition and the scope of the arguments}
\label{chap:current-scope}

The blockwise Alperin weight (BAW) conjecture predicts that a block has as many
irreducible Brauer characters as conjugacy classes of weights. The
inductive blockwise Alperin weight (iBAW) condition requires compatible
equivariant correspondences on
covering groups of a nonabelian simple group. These include conditions
on character extensions and on blocks of intermediate groups.

This report follows the manuscript's arguments for the inductive
blockwise Alperin weight condition for simple groups of Type B, simple
groups of Type C and sporadic simple groups. It begins with the general
preliminary results, then treats the three families and the reduction to
finite groups. The conclusions hold under the hypotheses stated in each
argument. References to manuscript results use the numbering of the accompanying
manuscript.
The mathematical library manual describes the foundations and the general
ModularRep results used in these arguments.

\section{Characters, weights and covering groups}

Fix a prime $\ell$. A weight of a finite group $G$ consists of a radical
$\ell$-subgroup $Q$ and an ordinary irreducible character $\theta$ of
$N_G(Q)/Q$ with defect zero. The set $\Alp(b)$ consists of the conjugacy
classes of weights belonging to a block $b$, and $\IBr(b)$ denotes its
irreducible Brauer characters. For passage through a quotient or an isomorphism, both the subgroup map and the character map must be specified under the stated hypotheses.

The full universal covering group of a simple group $S$ and its universal
$\ell'$-cover $X_\ell(S)$ play different roles. Each descent argument uses
the specified covering projection, central kernel and automorphism
action. The constructions for exceptional covers include these maps
and the character and block identifications they induce.

\begin{implementation}
The representations, irreducibility predicate and ordinary character
function are supplied by Mathlib. ModularRep uses these definitions to
form its type of ordinary irreducible character functions, and provides
the Brauer character and block constructions used below. The mathematical
library manual explains the imported definitions and results, the
additional constructions, and the explicit hypotheses.
\end{implementation}

\section{Coefficient fields and reduction}

Brauer characters use a chosen correspondence between roots of unity in
characteristic $\ell$ and characteristic zero. Character comparisons use
compatible roots.

The coefficient hypotheses depend on the application. The principal
$2$-block construction for Type C, including its rank two principal
case, uses an algebraically closed coefficient field $K$ for ordinary characters. It does not
require a discrete valuation ring with fraction field $K$. The applications
in defining characteristic and in low rank over even fields use a field
$K$ of characteristic zero, an algebraically closed field $k$ of
characteristic $\ell$, and the specified correspondences between roots
of unity. They do not require $K$ to be algebraically closed.
The application of Li--Li's theorem to all blocks in rank two and the applications
in higher rank at odd nondefining primes use a specified modular system
with fraction field $K$ and residue field $k$, together with the stated
splitting and root assumptions. Fixed embeddings of a common cyclotomic
field compare complex character values with values in $K$.

The hypotheses on the decomposition homomorphism, its naturality and
block induction refer to these coefficient choices. Where a published
statement uses complex characters, its application over another value
field includes the required character and root identifications.

\section{The conclusions and their assumptions}

The equality
\[
 |\IBr(b)|=|\Alp(b)|,
\]
an equivariant bijection between these sets, and a bijection satisfying
the prescribed block isomorphism relation between modular character
triples are distinct conclusions. The full inductive condition also requires compatible
extensions and block equalities for intermediate groups. At the trivial
radical subgroup it prescribes the correspondence coming from ordinary
characters of defect zero.

The arguments below state which conclusion is obtained. Counting
arguments give numerical equalities or equivariant bijections under the
stated group actions. Published criteria supply the further character
and block conditions under their hypotheses. When a passage from
individual blocks to a full family chooses compatible correspondences,
the resulting family need not coincide with independently chosen block
correspondences.

Each Lean theorem proves its conclusion under its stated hypotheses.
These include standard representation theory, cited results, finite
data and their interpretation for the groups in question. The geometric
constructions and trace identities used in manuscript Lemmas 2.11 and
2.12 remain assumptions. In particular, the construction of the centrally
normalised split pair in Type C is not formalised. Lean proves the finite spanning, independence,
coverage and character selection deductions. The conditional Type C
selection theorem concerns the specified family in the conformal
symplectic group. Induction and restriction identities give a separate
finite selection result for the symplectic group. The principal block
application takes the published field invariance and character counts
as explicit assumptions. The passage from the selected ordinary
characters to these principal block conclusions is not formalised.
At rank two, Li--Li's theorem is applied on the specified cover.
Compatibility across blocks and the quotient and root identifications
required for the central lift remain separate assumptions.
The inductive condition reported by CTBlocks
for the other $22$ sporadic groups is also assumed. A numerical table is
used with its stated interpretation and completeness assumptions.
The implementation paragraphs identify the deductions that are formalised.

\section{Scope of the hypotheses}

The Type C source interface applies the published block theorem only to
the specified automorphism stabilisers.
The theorem \lean{Main.theorem_1_1} assumes both
\lean{Main.ClassicalInputs} and \lean{Sporadic.Theorem.Inputs}.
The package does not construct either complete collection of source
assumptions. The conditional deductions do not establish the manuscript
without the stated external assumptions.

\section{References and source records}

Appendix~\ref{chap:current-statements} relates the numbered manuscript
statements to selected Lean declarations and states their roles.
Appendix~\ref{chap:current-source-index} indexes the manuscript citations
and bibliography. Appendix~\ref{app:construction-details} describes
further constructions in the source library and distinguishes auxiliary
results from those used in the manuscript deductions.
Numbers within Lean declaration names are not
manuscript cross-references. The source
records specify the version consulted, result location, relevant
hypotheses and use in the formalisation. A supporting reference
documents the stated result at that location. It is not, by itself,
a claim that the cited work first established the result.

\section{Compilation and verification}
\label{sec:compilation-resources}

All 1,999 project Lean source files and four audit sources were checked
with Lean 4.33.1 using \texttt{--trust=0}. The distributed Lean sources
agree with the compiled sources outside comments. The build reused
precompiled external dependencies at the pinned revisions. The recorded
external objects and direct imports were checked against their hashes.
The full tree of indirect external dependencies was not checked, and the
imported proofs were not replayed.

Compilation and verification used approximately 38.1 CPU hours. The peak
sampled sum of the processes' resident working sets was 10.40\,GiB.
Shared pages can be counted more than once, and peaks between samples
may be missed.

The package checker authenticates the distributed sources and metadata
without rerunning Lean, and the original build logs and compiled objects
are not distributed.

Source files for the supplementary exhaustive audit helper, its generator
and the generated audit module are not distributed. That audit cannot
be repeated using the supplied tooling.

For the audited declarations, the axiom checks restrict dependencies to
the permitted standard axioms. They do not check every declaration in the
library. The direct conclusion check rejects a leading hypothesis that
is definitionally equal to the conclusion.
It does not detect a merely logically equivalent hypothesis, a conclusion
hidden in a structure field, or joint circularity among hypotheses.
These checks do not establish the truth, independence or joint
satisfiability of the external assumptions.
